\documentclass[11pt,a4paper]{article}

% -------------------------------------------------
% Packages
% -------------------------------------------------
\usepackage[utf8]{inputenc}
\usepackage[T1]{fontenc}
\usepackage[brazil]{babel}
\usepackage[a4paper,top=1.2cm,bottom=1.2cm,left=1.4cm,right=1.4cm]{geometry}
\usepackage{hyperref}
\usepackage{xcolor}
\usepackage{titlesec}
\usepackage{enumitem}
\usepackage{tabularx}
\usepackage{fontawesome5}

\hypersetup{
    colorlinks=true,
    urlcolor=black
}

\definecolor{primary}{RGB}{25,55,109}

\pagestyle{empty}
\setlength{\parindent}{0pt}

\titleformat{\section}
{\Large\bfseries\color{primary}}
{}
{0pt}
{}
[\titlerule]

% Espaço antes/depois dos títulos de seção.
% Se passar de 2 páginas, diminua (ex.: {10pt}{5pt}); se sobrar espaço, aumente.
\titlespacing*{\section}{0pt}{18pt}{9pt}

\begin{document}

%%%%%%%%%%%%%%%%%%%%%%%%%%%%%%%%%%%%%%%%%%%%%%%%%%
% HEADER
%%%%%%%%%%%%%%%%%%%%%%%%%%%%%%%%%%%%%%%%%%%%%%%%%%

\begin{center}

{\Huge\bfseries Rafael Caetano da Silva}

\vspace{5pt}

{\Large Desenvolvedor .NET Pleno}

\vspace{4pt}

{\small
C\# • .NET Core • Angular • Blazor • SQL Server • Oracle • Docker • CI/CD
}

\vspace{6pt}

\small

\faEnvelope\
\href{mailto:rafaelkaetanodf@gmail.com}{rafaelkaetanodf@gmail.com}
\quad|\quad
\faPhone\
(61) 99398-9820
\quad|\quad
\faMapMarker*
Brasília - DF

\vspace{3pt}

\faGithub\
\href{https://github.com/Rafael-CaetanoS}{github.com/Rafael-CaetanoS}
\quad|\quad
\faLinkedin\
\href{https://linkedin.com/in/rafaelcaetanoo}{linkedin.com/in/rafaelcaetanoo}

\end{center}

\vspace{-0.2cm}
%%%%%%%%%%%%%%%%%%%%%%%%%%%%%%%%%%%%%%%%%%%%%%%%%%
\section{Resumo Profissional}

Desenvolvedor .NET Pleno com experiência em C\# .NET Core (ASP.NET Core), Angular, Blazor, Bootstrap, TypeScript e JavaScript no desenvolvimento de aplicações corporativas e integrações entre sistemas. Atuação na construção de APIs REST, bancos de dados SQL (Oracle, SQL Server, PostgreSQL e MySQL), mensageria com Azure Service Bus e pipelines de CI/CD com Docker e Azure DevOps. Recentemente, identifiquei e implementei uma solução de arquitetura que gerou economia de R\$ 6 mil mensais nos custos da empresa.

\vspace{0.15cm}

Sólida base em Orientação a Objetos e UML, Padrões de Projeto (Design Patterns), SOLID, Clean Architecture e Domain-Driven Design (DDD), com foco em Boas Práticas em Desenvolvimento de software, revisão de código e testes unitários com xUnit, além de Java com Spring Boot, JPA (Spring Data JPA) e Hibernate. Uso Agentes IA (Claude e ChatGPT) para ganhar produtividade no desenvolvimento. Participo da definição da arquitetura das soluções, do levantamento de requisitos técnicos e da tomada de decisões técnicas, exercendo atualmente funções de liderança técnica.

%%%%%%%%%%%%%%%%%%%%%%%%%%%%%%%%%%%%%%%%%%%%%%%%%%
\section{Experiência Profissional}

\textbf{Coordenador Back-end}
\hfill Ago/2025 -- Atual

Grupo Orion

\begin{itemize}[leftmargin=*,noitemsep,topsep=2pt]

\item \textbf{Economia de R\$ 6 mil mensais (\textasciitilde R\$ 72 mil/ano):} identifiquei e implementei uma solução para reduzir os custos da empresa, reestruturando a arquitetura dos sistemas, eliminando bancos de dados desnecessários e otimizando o DevOps.

\item Desenvolvimento e evolução de APIs REST em C\# .NET Core (ASP.NET Core) e de integrações entre sistemas corporativos, serviços externos e plataformas de monitoramento IoT.

\item Definição da arquitetura de novas funcionalidades utilizando Clean Architecture, Domain-Driven Design (DDD), SOLID e Padrões de Projeto (Design Patterns).

\item Estruturação e manutenção de pipelines de CI/CD (integração e entrega contínua) utilizando Azure DevOps.

\item Revisão de código, validação de soluções técnicas e promoção de boas práticas em desenvolvimento de software junto à equipe.

\item Tomada de decisões técnicas relacionadas à arquitetura, banco de dados, integrações, desempenho, custos e escalabilidade das aplicações.

\item Liderança técnica da equipe Back-end, com planejamento técnico das demandas e distribuição das atividades.

\item Levantamento e análise de requisitos técnicos junto às áreas de negócio e produto.

\end{itemize}

\vspace{0.2cm}

\textbf{Desenvolvedor Back-end Pleno}
\hfill Abr/2025 -- Ago/2025

Grupo Orion

\begin{itemize}[leftmargin=*,noitemsep,topsep=2pt]

\item Desenvolvimento de APIs REST e integrações entre sistemas com C\# .NET Core (ASP.NET Core), seguindo Clean Architecture e DDD.

\item Implementação de testes unitários utilizando xUnit para aumento da confiabilidade das aplicações.

\item Aplicação de Orientação a Objetos, SOLID, injeção de dependência e Padrões de Projeto (Design Patterns).

\item Modelagem e manutenção de bancos de dados SQL (PostgreSQL e MySQL).

\end{itemize}

\vspace{0.2cm}

\textbf{Analista de Sistemas (Estágio)}
\hfill Dez/2023 -- Dez/2024

Infraero

\begin{itemize}[leftmargin=*,noitemsep,topsep=2pt]

\item Desenvolvimento e manutenção de aplicações corporativas e sistemas internos de operações aeroportuárias com C\#, ASP.NET Core, Blazor e Java.

\item Desenvolvimento com bancos de dados Oracle, SQL Server, PostgreSQL e MySQL.

\item Utilização de Docker para padronização dos ambientes de desenvolvimento.

\item Versionamento de código utilizando Git e Azure DevOps.

\item Desenvolvimento de relatórios utilizando Microsoft Report Builder.

\end{itemize}

%%%%%%%%%%%%%%%%%%%%%%%%%%%%%%%%%%%%%%%%%%%%%%%%%%
\section{Projetos Relevantes}

\textbf{Redução de Custos de Infraestrutura -- Grupo Orion}

Identificação e implementação de uma solução para economizar nas contas da empresa: reestruturação da arquitetura dos sistemas, remoção de bancos de dados desnecessários e melhoria do DevOps, reduzindo os custos em R\$ 6 mil por mês.

\vspace{0.2cm}

\textbf{Automação de Climatização -- Mané Vírgula (Partage Shopping)}

Automação integrada à Plataforma Cortex, via WebSockets com SignalR, para controle em tempo real de 16 aparelhos de ar-condicionado do mercado Mané Vírgula, no Partage Shopping, pela aplicação ou por automações configuráveis pelos próprios usuários.

\vspace{0.2cm}

\textbf{Plataforma Cortex}

Responsável pelo desenvolvimento do backend da plataforma de monitoramento IoT, incluindo gerenciamento de alertas, notificações, configurações de alarmes e integração com dispositivos inteligentes.

\vspace{0.2cm}

\textbf{Integração Infraspeak}

Desenvolvimento de integrações entre sistemas via APIs REST para sincronização de ativos, ordens de serviço, falhas e propriedades adicionais entre sistemas corporativos.

\vspace{0.2cm}

\textbf{SIVWEB} -- Sistema de gerenciamento de monitores aeroportuários desenvolvido com ASP.NET Core e Blazor.

\vspace{0.2cm}

\textbf{SICOA} -- Sistema corporativo para identificação e controle de acesso de usuários.

%%%%%%%%%%%%%%%%%%%%%%%%%%%%%%%%%%%%%%%%%%%%%%%%%%
\section{Competências Técnicas}

\textbf{Back-end:} C\#, .NET Core, ASP.NET Core, APIs REST, Microsserviços, Java, Spring Boot

\vspace{0.12cm}

\textbf{Front-end:} Angular, Blazor, Bootstrap, TypeScript, JavaScript, HTML, CSS

\vspace{0.12cm}

\textbf{Arquitetura e Boas Práticas:} Orientação a Objetos, UML, Padrões de Projeto (Design Patterns), SOLID, Clean Architecture, Domain-Driven Design (DDD), CQRS, Arquitetura em Camadas, MVC

\vspace{0.12cm}

\textbf{Banco de Dados:} SQL, Oracle, SQL Server, PostgreSQL, MySQL

\vspace{0.12cm}

\textbf{ORM:} JPA (Spring Data JPA), Hibernate, Entity Framework, Dapper

\vspace{0.12cm}

\textbf{CI/CD e DevOps:} CI/CD (Integração e Entrega Contínua), Docker, Azure DevOps, Git, GitHub, Azure

\vspace{0.12cm}

\textbf{Testes:} Testes Unitários, xUnit

\vspace{0.12cm}

\textbf{Mensageria e Tempo Real:} Azure Service Bus, SignalR, WebSockets

\vspace{0.12cm}

\textbf{Ferramentas:} Agentes IA (Claude, ChatGPT), Swagger, Postman, Report Builder

%%%%%%%%%%%%%%%%%%%%%%%%%%%%%%%%%%%%%%%%%%%%%%%%%%
\section{Formação}

\textbf{Pós-graduação em Engenharia de Software}
\hfill 2025 -- 2026

Universidade Católica de Brasília

\vspace{0.2cm}

\textbf{Bacharelado em Sistemas de Informação}
\hfill 2021 -- 2024

Centro Universitário do Distrito Federal (UDF)

%%%%%%%%%%%%%%%%%%%%%%%%%%%%%%%%%%%%%%%%%%%%%%%%%%
\section{Idiomas}

Português — Nativo \quad|\quad Espanhol — Intermediário \quad|\quad Inglês — Básico

\end{document}
